\subsection{相对于交换因子的分次代数的张量积}

定义 6. 设 \((\Delta_{i})_{i\in I}\) 是一个以加法记号书写的交换幺半群的有限族。一个 \(\Delta_{i}\) 上的、取值于交换环 A 的交换因子系，是指一个映射系统 \(\varepsilon_{ij}:\Delta_i\times \Delta_j\to \mathbf{A}\) （其中 \(i\in \mathbf{I}, j\in \mathbf{I}, i\neq j\) ），满足以下条件：

\[
\varepsilon_ {i j} \left(\alpha_ {i} + \alpha_ {i} ^ {\prime}, \beta_ {j}\right) = \varepsilon_ {i j} \left(\alpha_ {i}, \beta_ {j}\right) \varepsilon_ {i j} \left(\alpha_ {i} ^ {\prime}, \beta_ {j}\right)\tag{22}
\]

\[
\varepsilon_ {i j} \left(\alpha_ {i}, \beta_ {j} + \beta_ {j} ^ {\prime}\right) = \varepsilon_ {i j} \left(\alpha_ {i}, \beta_ {j}\right) \varepsilon_ {i j} \left(\alpha_ {i}, \beta_ {j} ^ {\prime}\right)\tag{23}
\]

\[
\varepsilon_ {i j} (\alpha_ {i}, \beta_ {j}) \varepsilon_ {j i} (\beta_ {j}, \alpha_ {i}) = 1,\tag{24}
\]

对所有 \(\alpha_{i},\alpha_{i}^{\prime}\) （属于 \(\Delta_{i},\beta_{j},\beta_{j}^{\prime}\) ）以及 \(\Delta_j\) 中的元素成立。

若给 I 赋予一个全序且 \(\Delta_{i}\) 是群，则在 \(\Delta_{i}\) 上可以如下定义一个交换因子系：对每个有序对 \((i,j)\) （满足 i\textless j ），任取一个从 \(\Delta_{i}\times\Delta_{j}\) 到由环 A 的可逆元素组成的（乘法）Z-模 A* 的 Z-双线性映射，然后写出

\[
\varepsilon_ {j i} (\beta_ {j}, \alpha_ {i}) = (\varepsilon_ {i j} (\alpha_ {i}, \beta_ {j})) ^ {- 1}
\]

对 \(i < j\) .

注意，由于 \(\varepsilon_{ij}(\alpha_{i},\beta_{j})\) 都是可逆的，

\[
\varepsilon_ {i j} (0, \beta_ {j}) = \varepsilon_ {i j} (\alpha_ {i}, 0) = 1,
\]

由 (22) 与 (23) 即得。

例。(1) 平凡的交换因子系由满足下列条件的 \(\varepsilon_{ij}\) 组成： \(\varepsilon_{ij}(\alpha_i, \beta_j) = 1\) 对所有 \(i, j, \alpha_i \in \Delta_i, \beta_j \in \Delta_j\) 成立.

\begin{enumerate}
\def\labelenumi{(\arabic{enumi})}
\setcounter{enumi}{1}
\tightlist
\item
  若取 \(\mathbf{A} = \mathbf{Z}\) 且取 \(\Delta_{i} = \mathbf{Z}\) （对所有 \(i \in \mathbf{I}\) ），则令 \(\varepsilon_{ij}(\alpha_i, \beta_j) = (-1)^{\alpha_i\beta_j}\) 即得到一个交换因子系。注意，这个数仅依赖于 \(\alpha_{i}\) 和 \(\beta_{j}\) 的奇偶性，因此 \(\varepsilon_{ij}\) 在某些 \(\Delta_{i}\) 等于 \(\mathbf{Z}/2\mathbf{Z}\) 而其余的等于 \(\mathbf{Z}\) 的情形下仍可作为交换因子。
\end{enumerate}

这两个例子是应用中最常见的情况。

命题 10. 设 A 是一个交换环， \((\Delta_{i})_{i\in I}\) 是一个以加法记号书写的交换幺半群的有限族；对每个 \(i\in I\) ，设 \(\mathbf{E}_i\) 是一个类型为 \(\Delta_i\) 的分次 A-代数。最后，设 \((\varepsilon_{ij})\) 是 \(\Delta_i\) 上的一个取值于 A 的交换因子系。则存在一个类型为 \(\Delta = \prod_{i\in I}\Delta_i\) 的分次 A-代数 E，以及对每个 \(i\in I\) 的一个代数同态 \(h_i:\mathbf{E}_i\to \mathbf{E}\) ，具有下列性质：

\begin{enumerate}
\def\labelenumi{(\roman{enumi})}
\item
  若 \(\phi_{i}:\Delta_{i}\to\Delta\) 是典范同态，则 \(h_{i}\) 是分次同态（II, § 11, 第 2 号），换言之， \(h_{i}(E_{i}^{\alpha_{i}})\subset E^{\phi_{i}(\alpha_{i})}\) ，其中 \((E_{i}^{\alpha_{i}})\) 与 \((E^{\alpha})\) 分别表示 \(E_{i}\) 与 \(E\) 上的分次。
\item
  若 \(i \neq j\) 且 \(x_{i}\) （相应地 \(x_{j}\) ）是 \(\mathbf{E}_i\) （相应地 \(\mathbf{E}_j\) ）的次数为 \(\alpha_i \in \Delta_i\) （相应地 \(\beta_j \in \Delta_j\) ）的齐次元素，则
\end{enumerate}

\[
h _ {i} (x _ {i}) h _ {j} (x _ {j}) = \varepsilon_ {i j} (\alpha_ {i}, \beta_ {j}) h _ {j} (x _ {j}) h _ {i} (x _ {i}).\tag{25}
\]

\begin{enumerate}
\def\labelenumi{(\roman{enumi})}
\setcounter{enumi}{2}
\tightlist
\item
  对每个 A-代数 \(\mathbf{F}\) 以及每个满足下列条件的同态系统 \(f_{i}:\mathbf{E}_{i}\to \mathbf{F}\)
\end{enumerate}

\[
f _ {i} (x _ {i}) f _ {j} (x _ {j}) = \varepsilon_ {i j} (\alpha_ {i}, \beta_ {j}) f _ {j} (x _ {j}) f _ {i} (x _ {i}),\tag{26}
\]

其中 \(i, j, x_i, x_j, \alpha_i, \beta_j\) 如 (ii) 中所述，则存在且仅存在一个代数同态 \(f: \mathbf{E} \to \mathbf{F}\) ，使得 \(f_i = f \circ h_i\) 对所有 \(i \in \mathbf{I}\) 成立。此外， \(\mathbf{E}\) 的底 A-模是张量积 \(\bigotimes_{i \in \mathbf{I}} \mathbf{E}_i\) .

考虑 A-模 \(E = \bigotimes_{i \in I} E_i\) ；它等同于子模 \(E^\alpha\) 的直和，其中对每个 \(\alpha = (\alpha_i) \in \Delta\) ，记 \(E^\alpha = \bigotimes_{i \in I} E_i^{\alpha_i}\) ；因而这些 \(E^\alpha\) 在 A-模 E 上构成一个类型为 \(\Delta\) 的分次。我们将要在 E 上定义一个类型为 \(\Delta\) 的分次 A-代数结构。为此，给 I 一个全序；对每一对元素 \(\alpha = (\alpha_i)\) , \(\beta = (\beta_i)\) （属于 \(\Delta\) ），我们必须首先定义一个从 \(E^\alpha \times E^\beta\) 到 \(E^\alpha + \beta\) 的 A-双线性映射，或者等价地，定义一个 A-线性映射 \(m_{\alpha\beta}\) ，它把 \(E^\alpha \otimes_A E^\beta\) 映到 \(E^\alpha + \beta\) 。我们将用条件来定义 \(m_{\alpha\beta}\)

\[
m _ {\alpha \beta} \left(\left(\bigotimes_ {i \in I} x _ {i}\right) \otimes \left(\bigotimes_ {i \in I} y _ {i}\right)\right) = \varepsilon (\alpha , \beta) \bigotimes_ {i \in I} (x _ {i} y _ {i})\tag{27}
\]

对于 \(x_{i}\in \mathrm{E}_{i}^{\alpha_{i}},y_{i}\in \mathrm{E}_{i}^{\beta_{i}}\) ，其中

\[
\varepsilon (\alpha , \beta) = \prod_ {i > j} \varepsilon_ {i j} (\alpha_ {i}, \beta_ {j}).\tag{28}
\]

(27) 的右边显然属于 \(\mathbf{E}^{\alpha +\beta}\) ，且映射 \((x_{1},\ldots ,x_{n},y_{1},\ldots ,y_{n})\mapsto \varepsilon (\alpha ,\beta)\bigotimes_{i\in I}(x_{i}y_{i})\) 在诸 \(\mathrm{E}_i^{\alpha_i}\) 与诸 \(\mathrm{E}_i^{\beta_i}(1\leqslant i\leqslant n)\) 的乘积中是 A-多重线性的。接着必须证明这样定义在 E 上的乘法是结合的；现在，若 \(\gamma = (\gamma_{i})\) 是 \(\Delta\) 的第三个元素，且 \(z_{i} \in \mathrm{E}_{i}^{\gamma_{i}}\) （对 \(1 \leqslant i \leqslant n\) ），则

\[
\left(\left(\bigotimes_ {i} x _ {i}\right) \left(\bigotimes_ {i} y _ {i}\right)\right) \left(\bigotimes_ {i} z _ {i}\right) = \varepsilon (\alpha + \beta , \gamma) \varepsilon (\alpha , \beta) \bigotimes_ {i} (x _ {i} y _ {i} z _ {i})
\]

\[
\left(\bigotimes_ {i} x _ {i}\right) \left(\left(\bigotimes_ {i} y _ {i}\right) \left(\bigotimes_ {i} z _ {i}\right)\right) = \varepsilon (\alpha , \beta + \gamma) \varepsilon (\beta , \gamma) \bigotimes_ {i} (x _ {i} y _ {i} z _ {i})
\]

从而只需验证恒等式

\[
\varepsilon (\alpha + \beta , \gamma) \varepsilon (\alpha , \beta) = \varepsilon (\alpha , \beta + \gamma) \varepsilon (\beta , \gamma).
\]

但后者立即由以下关系得出

\[
\varepsilon (\alpha + \beta , \gamma) = \varepsilon (\alpha , \beta) \varepsilon (\beta , \gamma)
\]

\[
\varepsilon (\alpha , \beta + \gamma) = \varepsilon (\alpha , \beta) \varepsilon (\alpha , \gamma)
\]

这些关系本身是定义(28)、(22)和(23)的直接推论。

若对所有 \(i \in \mathbf{I}\) ， \(e_i\) 表示 \(\mathbf{E}_i\) 的单位元，我们知道 \(e_i\) 是 0 次齐次的（§ 3, 第 1 号），因此 \(e = \bigotimes_{i \in \mathbf{I}} e_i\) 是 0 次齐次的，并且由 (27)、(28) 及关系式

\[
\varepsilon_ {i j} (\alpha_ {i}, 0) = \varepsilon_ {i j} (0, \beta_ {j}) = 1
\]

即 \(e\) 是 \(\mathbf{E}\) 的单位元，这就完成了 \(\mathbf{E}\) 上所需要的分次 A-代数结构的定义。然后取 \(h_i(x_i) = x_i \otimes \bigotimes_{j \neq i} e_j\) ；为验证 \(h_i(x_i x_i') = h_i(x_i) h_i(x_i')\) 对 \(x_i, x_i'\) （属于 \(\mathbf{E}_i\) ）成立，只需限于 \(x_i\) 和 \(x_i'\) 均为齐次元素的情形，此时该关系由 (27) 和关系式 \(\varepsilon_{ij}(\alpha_i, 0) = \varepsilon_{ij}(0, \beta_j) = 1\) 立即得出；同样的关系式以及 (24) 还证明了 \(h_i\) 满足命题中的条件 (i) 与 (ii)，且

\[
\bigotimes_ {i \in I} x _ {i} = \prod_ {i \in I} h _ {i} (x _ {i})\tag{29}
\]

其中右端是有序序列 \((h_i(x_i))_{i\in I}\) 按 I 上给定的全序在 E 中取的乘积（I, § 1, 第 2 号）（只需对 \(x_{i}\) （设为齐次）中异于 \(e_i\) 者的个数作归纳论证即可）.

仍需证明条件 (iii)；注意，映射

\[
(x _ {i}) _ {i \in I} \mapsto \prod_ {i \in I} f _ {i} (x _ {i}),
\]

其中右端是有序序列 \((f_{i}(x_{i}))_{i\in\mathbf{I}}\) 按 I 上给定的全序取的乘积，是 A-多重线性的。于是存在且仅存在一个 A-线性映射 \(f: E \rightarrow F\) ，使得

\[
f \left(\bigotimes_ {i \in I} x _ {i}\right) = \prod_ {i \in I} f _ {i} (x _ {i}).\tag{30}
\]

显然 \(f(e)\) 是 \(\mathbf{F}\) 的单位元，且 \(f \circ h_i = f_i\) ；为验证 \(f\) 是代数同态，换言之，验证 \(f(x)f(y) = f(xy)\) 对 \(x, y\) （属于 \(\mathbf{E}\) ）成立，由线性性，只需限于下列情形： \(x = \bigotimes_{i \in \mathbf{I}} x_i\) 且 \(y = \bigotimes_{i \in \mathbf{I}} y_i\) ，其中 \(x_i\) （相应地 \(y_i\) ）是次数为 \(\alpha_i\) （相应地 \(\beta_i\) ）的齐次元素（对所有 \(i \in \mathbf{I}\) ）。此时，要验证的关系式由 (27) 化归为

\[
\left(\prod_ {i \in \mathrm{I}} f _ {i} (x _ {i})\right) \left(\prod_ {i \in \mathrm{I}} f _ {i} (y _ {i})\right) = \varepsilon (\alpha , \beta) \prod_ {i \in \mathrm{I}} \left(f _ {i} (x _ {i}) f _ {i} (y _ {i})\right).
\]

但考虑到关系式 (26)，这是第 6 号的引理 2 的一个推论。

显然，代数 E 与典范映射 \(\psi: \bigotimes_{i \in I} E_i \to E\) 构成泛映射问题（集合论, IV, § 3, 第 1 号）的一个解，其中 \(\Sigma\) 是 A-代数结构的种类，而 \(\alpha\) -映射是映射 \(\prod_{i} f_i\) （从 \(\prod_{i} E_i\) 到一个 A-代数，满足条件 (26)）。

对 I 上固定的一个全序，命题 10 的证明中所定义的分次代数 E 将称为一个分次张量 \(\varepsilon\) -积，其类型为 \(\Delta\) ，所属的族为 \((\mathrm{E}_i)_{i\in \mathbf{I}}\) （族中各代数的类型为 \(\Delta_i\) ），并记作 \({}^{\varepsilon}\bigotimes_{i\in \mathbf{I}}\mathrm{E}_i\) （若关于 I 上的序不致引起混淆）；类似地，命题 10 的证明中定义的同态 \(f:\mathrm{E}\to \mathrm{F}\) 将记作 \({}^{\varepsilon}\bigotimes_{i\in \mathbf{I}}f_i\) 。同态 \(h_i\) 称为典范的。我们也写作 \({}^{\varepsilon}\mathrm{G}^{\otimes n}\) （当 \(\mathbf{I} = [1,n]\) 且所有 \(\mathrm{E}_i\) 都等于同一个代数 G 时）。

注。(1) 若取 \(\varepsilon_{ij}(\alpha_i, \beta_j) = 1\) （对所有 \(i, j, \alpha_i\) 与 \(\beta_j\) ），我们就重新得到第 1 号中定义的代数的张量积（其分次还是各因子的分次的张量积）.

\begin{enumerate}
\def\labelenumi{(\arabic{enumi})}
\setcounter{enumi}{1}
\tightlist
\item
  设所有 \(\Delta_{i}\) 都等于 \(\mathbf{Z}\) ，并写 \(\varepsilon_{ij}(\alpha_i, \beta_j) = (-1)^{\alpha_i \beta_j}\) ；对应于这个交换因子系的张量 \(\varepsilon\) -积 \({}^{\varepsilon}\bigotimes_{i \in I} E_i\) 称为分次代数 \(E_i\) （类型为 \(\mathbf{Z}\) ）的斜张量积，记作 \({}^{g}\bigotimes_{i \in I} E_i\) （对两个代数则记 \(E{}^{g} \otimes_A F\) ，或以 \({}^{g} G^{\otimes n}\) 代替 \({}^{\varepsilon} G^{\otimes n}\) ）.
\end{enumerate}

推论 1. 在命题 10 的记号下，进一步设 F 是一个类型为 \(\Delta\) 的分次 A-代数，且每个 \(f_{i}\) 都是相对于 \(\phi_{i}:\Delta_{i}\to\Delta\) 的分次代数同态；则 \(f={}^{\varepsilon}\bigotimes_{i}f_{i}\) 是一个分次代数同态。

这由 \(f\) 的定义以及下列事实立即得出：

\[
\sum_ {i \in I} \phi_ {i} (\alpha_ {i}) = (\alpha_ {i})
\]

这是由 \(\phi_{i}\) 的定义得到的.

由此可以看出， \((\mathrm{E},\psi)\) 也是另一个泛映射问题的解，这一次 \(\Sigma\) 是类型为 \(\Delta\) 的分次 A-代数结构的种类，其态射是类型为 \(\Delta\) 的分次代数同态，而 \(\alpha\) -映射是映射 \(\prod_{i}f_{i}\) （其中除条件 (26) 外，还假设 \(f_{i}\) 是相对于 \(\phi_{i}\) 的分次代数同态）.

推论 2. 设 \((\mathbf{E}_i)_{i\in \mathbf{I}}\) , \((\mathbf{F}_i)_{i\in \mathbf{I}}\) 是两个有限的 A-代数族，且 \(\mathbf{E}_i\) 与 \(\mathbf{F}_i\) 都是类型为 \(\Delta_i\) 的分次代数（对所有 \(i\in \mathbf{I}\) ）。对每个 \(i\in \mathbf{I}\) ，设 \(g_i:\mathbf{E}_i\to \mathbf{F}_i\) 是一个类型为 \(\Delta_i\) 的分次代数同态。那么，若 \(h_i:\mathbf{E}_i\to {}^{\varepsilon}\bigotimes_{i\in \mathbf{I}}\mathbf{E}_i\) 与 \(h_i':\mathbf{F}_i\to {}^{\varepsilon}\bigotimes_{i\in \mathbf{I}}\mathbf{F}_i\) 是典范同态，则存在且仅存在一个类型为 \(\Delta\) 的分次 A-代数同态 \(g: {}^{\varepsilon}\bigotimes_{i\in \mathbf{I}}\mathbf{E}_i\to {}^{\varepsilon}\bigotimes_{i\in \mathbf{I}}\mathbf{F}_i\) ，使得 \(g\circ h_i = h_i'\circ g_i\) 对所有 \(i\in \mathbf{I}\) 成立。此外，若每个 \(g_i\) 都是双射，则 \(g\) 也是.

只需将推论 1 应用于 \(f_{i} = h_{i}^{\prime} \circ g_{i}\) ，并注意到条件 (26) 可由应用于 \(h_{i}^{\prime}\) 的关系式 (25) 得出。

推论 2 中定义的同态也记作 \({}^{\varepsilon}\bigotimes_{i} g_{i}\) （若不致引起混淆）；若对每个 \(i \in I\) ， \(G_{i}\) 是又一个分次 A-代数（类型为 \(\Delta_{i}\) ），且 \(g_{i}^{\prime}: F_{i} \to G_{i}\) 是一个分次代数同态，则

\[
\left(^ {\varepsilon} \bigotimes_ {i} g _ {i} ^ {\prime}\right) \circ \left(^ {\varepsilon} \bigotimes_ {i} g _ {i}\right) = ^ {\varepsilon} \bigotimes_ {i} \left(g _ {i} ^ {\prime} \circ g _ {i}\right),\tag{31}
\]

这由(30)立即得出。

在斜张量积（类型为 \(\mathbf{Z}\) 的分次代数的）情形，我们写 \({}^{g}\bigotimes_{i}f_{i}\) 以代替 \({}^{\varepsilon}\bigotimes_{i}f_{i}\) ——对同态 \(f_{i}: \mathrm{E}_{i} \to \mathrm{F}_{i}\) （类型为 \(\mathbf{Z}\) 的分次代数的同态）而言；当 \(\mathrm{I} = \{1, 2\}\) 时，这个同态也记作 \(f_{1}{}^{g}\otimes f_{2}\) ；当 \(\mathrm{I} = [1, n]\) 且所有 \(\mathrm{E}_{i}\) （相应地 \(\mathrm{F}_{i}\) ）相等并且所有 \(f_{i}\) 都等于同一个同态 \(f\) 时，我们写 \({}^{g}f^{\otimes n}\) .

注。在命题 10 的证明中，I 上的一个全序被用来在张量积 \(\bigotimes_{i\in I}E_i\) 上定义一个代数结构，该张量积是 A-模

\(E_{i}\) 的张量积。若改变 I 上的序，则在 \(\bigotimes_{i\in I}E_{i}\) 上会产生另一个乘法结构，但这样得到的新代数与上述代数典范同构，因为二者都是同一个泛映射问题的解。例如，当 \(I=\{1,2\}\) 时，从代数 \(E_{1}{}^{\varepsilon}\otimes_{A}E_{2}\) 到代数 \(E_{2}{}^{\varepsilon}\otimes_{A}E_{1}\) 的典范同构把 \(x_{1}\otimes x_{2}\) 映为 \(\varepsilon_{2,1}(\alpha,\beta)x_{2}\otimes x_{1}\) ，其中 \(x_{1}\) 是次数为 \(\alpha\) 的齐次元素， \(x_{2}\) 是次数为 \(\beta\) 的齐次元素.

设 J 是 I 的一个子集，且对每个 \(i \in J\) ，考虑典范同态 \(h_{i}: E_{i} \to {}^{\varepsilon}\bigotimes_{i \in I} E_{i} = E\) 。借助关系式 (25)，从这些同态典范地（由命题 10）导出一个典范同态 \(h: E' = {}^{\varepsilon}\bigotimes_{i \in J} E_{i} \to E\) ，使得对所有 \(i \in J\) ， \(h_{i}' = h \circ h_{i}\) ，其中 \(h_{i}'\) 是典范同态 \(E_{i} \to E'\) 。取 J 上由 I 上所选的全序诱导的全序，我们得到

\[
h \left(\bigotimes_ {i \in J} x _ {i}\right) = \prod_ {i \in I} h _ {i} \left(x _ {i}\right) = \bigotimes_ {i \in I} x _ {i} ^ {\prime}
\]

其中中间项是有序序列 \((h_i(x_i))_{i\in J}\) 的乘积，而在右端项中， \(x_{i}^{\prime} = x_{i}\) 对 \(i\in J\) 成立， \(x_{i}^{\prime} = e_{i}\) 对 \(i\notin J\) 成立.

命题 11.（张量 \(\varepsilon\) -积的“结合性”）。在命题 10 的记号下，设 \((J_{\lambda})_{\lambda \in L}\) 是 I 的一个划分，并记 \(\Delta_{\lambda}' = \prod_{i \in J_{\lambda}} \Delta_i\) （对所有 \(\lambda \in L\) ）。设 \(E_{\lambda}'\) 是一个分次张量 \(\varepsilon\) -积，其类型为 \(\Delta_{\lambda}'\) ，所属的族为 \((E_i)_{i \in J_{\lambda}}\) （对 \(J_{\lambda}\) 上选定的某个全序而言）。另一方面，对 L 中的 \(\lambda, \mu\) 且 \(\lambda \neq \mu\) ，我们对 \(\alpha_{\lambda}' = (\alpha_i)_{i \in J_{\lambda}}, \beta_{\mu}' = (\beta_j)_{j \in J_{\mu}}\) 写

\[
\varepsilon_ {\lambda \mu} ^ {\prime} (\alpha _ {\lambda} ^ {\prime}, \beta_ {\mu} ^ {\prime}) = \prod_ {i \in J _ {\lambda}, j \in J _ {\mu}} \varepsilon_ {i j} (\alpha_ {i}, \beta_ {j}).\tag{32}
\]

那么 \((\varepsilon_{\lambda \mu}^{\prime})\) 是 \(\Delta_{\lambda}^{\prime}\) 上的一个取值于 A 的交换因子系，且存在且仅存在一个类型为 \(\Delta, v: {}^{\varepsilon'}\bigotimes_{\lambda \in L} E_{\lambda}^{\prime} \to {}^{\varepsilon}\bigotimes_{i \in I} E_{i}\) 的分次代数同态，使得

\[
v \left(\bigotimes_ {\lambda \in L} \left(\bigotimes_ {i \in J _ {\lambda}} x _ {i}\right)\right) = \bigotimes_ {i \in I} x _ {i}\tag{33}
\]

对所有 \((x_{i})\in \prod_{i\in \mathbf{I}}\mathrm{E}_{i}\) 成立，前提是在 \(\mathbf{I}\) 上取这样的全序：它在每个 \(\mathbf{J}_{\lambda}\) 上诱导所选的全序，并且使得对 \(\lambda < \mu\) （在 \(\mathbf{L}\) 中）、 \(i\in \mathbf{J}_{\lambda}\) 与 \(j\in \mathbf{J}_{\mu}\) ，有 \(i < j\) .

\(\varepsilon_{\lambda \mu}'\) 构成交换因子系这一事实是平凡的。设 \(h_{i,\lambda}:\mathrm{E}_i\to \mathrm{E}_{\lambda}', h_{\lambda}':\mathrm{E}_{\lambda}'\to {}^{\varepsilon'}\bigotimes_{\lambda \in L}\mathrm{E}_{\lambda}'\) 是典范同态（对 \(\lambda \in L\) , \(i\in \mathbf{J}_{\lambda}\) ），并记 \(h_i'' = h_\lambda' \circ h_{i,\lambda}\) ；由泛映射问题解的唯一性，只需证明 \({}^{\varepsilon'}\bigotimes_{\lambda \in L} E_{\lambda}'\) 与诸 \(h_i''\) 满足命题 10 的条件即可。现在，对所有 \(\lambda \in L\) ，设 \(f_{\lambda}': E_{\lambda}' \to F\) 是唯一的代数同态，使得 \(f_{\lambda}' \circ h_{i,\lambda} = f_i\) 对所有 \(i \in J_{\lambda}\) 成立。我们证明，对 \(\lambda \neq \mu\) , \(\alpha_{\lambda}' = (\alpha_i)_{i \in J_{\lambda}}\) , \(\beta_{\mu}' = (\beta_j)_{j \in J_{\mu}}\) ,

\[
f _ {\lambda} ^ {\prime} (x _ {\lambda} ^ {\prime}) f _ {\mu} ^ {\prime} (x _ {\mu} ^ {\prime}) = \varepsilon_ {\lambda \mu} ^ {\prime} (\alpha_ {\lambda} ^ {\prime}, \beta_ {\mu} ^ {\prime}) f _ {\mu} ^ {\prime} (x _ {\mu} ^ {\prime}) f _ {\lambda} ^ {\prime} (x _ {\lambda} ^ {\prime})\tag{34}
\]

对齐次元素 \(x_{\lambda}^{\prime} \in \mathrm{E}_{\lambda}^{\prime}\) （相应地 \(x_{\mu}^{\prime} \in \mathrm{E}_{\mu}^{\prime}\) ），其次数为 \(\alpha_{\lambda}^{\prime}\) （相应地 \(\beta_{\mu}^{\prime}\) ）；由线性性，只需在 \(x_{\lambda}^{\prime} = \bigotimes_{i \in J_{\lambda}} x_{i}\) , \(x_{\mu}^{\prime} = \bigotimes_{j \in J_{\mu}} x_{j}\) 的情形验证，其中 \(x_{i}\) （相应地 \(x_{j}\) ）是次数为 \(\alpha_{i}\) （相应地 \(\beta_{j}\) ）的齐次元素，属于 \(E_{i}\) （相应地 \(E_{j}\) ）（对 \(i \in J_{\lambda}, j \in J_{\mu}\) ）。但这由定义诸 \(f_{\lambda}^{\prime}\) 的公式 (30) 以及第 6 号的引理 3 得出，其中用到假设 (26) 与定义 (32)。因此存在且仅存在一个代数同态 \(f: ^{\varepsilon'} \bigotimes_{\lambda \in L} E_{\lambda}' \to F\) ，使得 \(f \circ h_{\lambda}' = f_{\lambda}'\) 对所有 \(\lambda \in L\) 成立；由此 \(f \circ h_{i}' = f_{i}\) 对所有 \(i \in I\) 成立，而 \(f\) 的唯一性是平凡的。

